% GENERATED FILE. Edit canonical lessons, then run npm run generate:books.
% canonical-source-hash: fnv1a64:c5509e334219cb33
% canonical-lessons: GE-C286-urteil, GE-C286-richter, GE-C286-verurteilen, GE-C286-schuldig, GE-C286-unschuldig

\chapter{Vor Gericht}
\label{ch:ge-vor-gericht}
\hlchaptermodality{286}

\begin{chapteropening}
\textbf{By the end of this chapter:} \emph{I can report a trial: a verdict, a judge, to sentence, guilty and innocent.}

Complete the last lesson of chapter 286: I can report a trial: a verdict, a judge, to sentence, guilty and innocent.
\end{chapteropening}

\section[das Urteil]{das Urteil — the verdict, the judgement}

\label{lesson:GE-C286-urteil}



\begin{quote}
Before the new one: say the German for to forge, then the German for to blackmail.
\end{quote}

\subsection*{You'll want to know: das Urteil}
\textbf{das Urteil} — \textquotedblleft{}the verdict, the judgement\textquotedblright{}.

\begin{quote}
Das Urteil kam am Freitag.
\end{quote}

\begin{quote}
\emph{The verdict came on Friday.}
\end{quote}

Plural: \emph{die Urteile}. \emph{Ein Urteil fällen} — to pass judgement.

\begin{cousinweb}
\emph{Urteil} goes back to Proto-Germanic \emph{uzdailiją} \textquotedblleft{}judgement\textquotedblright{} (EtymDB), from \emph{uz-} \textquotedblleft{}out\textquotedblright{} and the root of \emph{Teil} \textquotedblleft{}part\textquotedblright{}: what is dealt out. English \emph{ordeal} is a cousin.
\end{cousinweb}

\begin{grammarlens}[title={What you've built}]
One more word for telling, explaining and describing at B1.
\end{grammarlens}

\subsection*{Your turn}
\begin{itemize}
\raggedright
  \item \emph{Say it:} \emph{das Urteil}
  \item \emph{Say it:} \emph{das Urteil}, once more
  \item \emph{Say it:} a sentence of your own with \emph{das Urteil}
  \item \emph{Recall:} say the German for to forge, then the German for to blackmail, then say \emph{das Urteil} again
\end{itemize}

\begin{tcolorbox}[breakable,colback=teal!4,colframe=teal!35!black,title={Before you move on}]
What does das Urteil mean? (\textquotedblleft{}The verdict, the judgement\textquotedblright{}.) Say it once more.
\end{tcolorbox}
\section[der Richter]{der Richter — the judge}

\label{lesson:GE-C286-richter}



\begin{quote}
Before the new one: say the German for to blackmail, then the German for the verdict.
\end{quote}

\subsection*{You'll want to know: der Richter}
\textbf{der Richter} — \textquotedblleft{}the judge\textquotedblright{}.

\begin{quote}
Der Richter hörte beiden Seiten zu.
\end{quote}

\begin{quote}
\emph{The judge listened to both sides.}
\end{quote}

Plural: \emph{die Richter}. Feminine \emph{die Richterin}.

\begin{cousinweb}
\emph{Richter} is built on \emph{richten} \textquotedblleft{}to judge, to put right\textquotedblright{}, from Old High German \emph{rihten}, Proto-Germanic \emph{rihtijaną} (EtymDB).
\end{cousinweb}

\begin{grammarlens}[title={What you've built}]
One more word for telling, explaining and describing at B1.
\end{grammarlens}

\subsection*{Your turn}
\begin{itemize}
\raggedright
  \item \emph{Say it:} \emph{der Richter}
  \item \emph{Say it:} \emph{der Richter}, once more
  \item \emph{Say it:} a sentence of your own with \emph{der Richter}
  \item \emph{Recall:} say the German for to blackmail, then the German for the verdict, then say \emph{der Richter} again
\end{itemize}

\begin{tcolorbox}[breakable,colback=teal!4,colframe=teal!35!black,title={Before you move on}]
What does der Richter mean? (\textquotedblleft{}The judge\textquotedblright{}.) Say it once more.
\end{tcolorbox}
\section[verurteilen]{verurteilen — to sentence, to condemn}

\label{lesson:GE-C286-verurteilen}



\begin{quote}
Before the new one: say the German for the verdict, then the German for the judge.
\end{quote}

\subsection*{You'll want to know: verurteilen}
\textbf{verurteilen} — \textquotedblleft{}to sentence, to condemn\textquotedblright{}.

\begin{quote}
Der Richter verurteilte ihn zu einer hohen Strafe.
\end{quote}

\begin{quote}
\emph{The judge sentenced him to a heavy penalty.}
\end{quote}

\emph{Zu etwas verurteilen} — to sentence to something.

\begin{cousinweb}
\emph{verurteilen} is built on \emph{Urteil} \textquotedblleft{}verdict\textquotedblright{}, from Proto-Germanic \emph{uzdailiją} \textquotedblleft{}judgement\textquotedblright{} (EtymDB), with the prefix \emph{ver-}.
\end{cousinweb}

\begin{grammarlens}[title={What you've built}]
One more word for telling, explaining and describing at B1.
\end{grammarlens}

\subsection*{Your turn}
\begin{itemize}
\raggedright
  \item \emph{Say it:} \emph{verurteilen}
  \item \emph{Say it:} \emph{verurteilen}, once more
  \item \emph{Say it:} a sentence of your own with \emph{verurteilen}
  \item \emph{Recall:} say the German for the verdict, then the German for the judge, then say \emph{verurteilen} again
\end{itemize}

\begin{tcolorbox}[breakable,colback=teal!4,colframe=teal!35!black,title={Before you move on}]
What does verurteilen mean? (\textquotedblleft{}To sentence, to condemn\textquotedblright{}.) Say it once more.
\end{tcolorbox}
\section[schuldig]{schuldig — guilty; owing}

\label{lesson:GE-C286-schuldig}



\begin{quote}
Before the new one: say the German for the judge, then the German for to sentence.
\end{quote}

\subsection*{You'll want to know: schuldig}
\textbf{schuldig} — \textquotedblleft{}guilty; owing\textquotedblright{}.

\begin{quote}
Das Gericht fand den Mann schuldig.
\end{quote}

\begin{quote}
\emph{The court found the man guilty.}
\end{quote}

Also 'owing': \emph{Du bist mir zehn Euro schuldig}.

\begin{cousinweb}
\emph{schuldig} is inherited from Old High German \emph{sculdig} (EtymDB), built on \emph{Schuld} \textquotedblleft{}guilt, debt\textquotedblright{}.
\end{cousinweb}

\begin{grammarlens}[title={What you've built}]
One more word for telling, explaining and describing at B1.
\end{grammarlens}

\subsection*{Your turn}
\begin{itemize}
\raggedright
  \item \emph{Say it:} \emph{schuldig}
  \item \emph{Say it:} \emph{schuldig}, once more
  \item \emph{Say it:} a sentence of your own with \emph{schuldig}
  \item \emph{Recall:} say the German for the judge, then the German for to sentence, then say \emph{schuldig} again
\end{itemize}

\begin{tcolorbox}[breakable,colback=teal!4,colframe=teal!35!black,title={Before you move on}]
What does schuldig mean? (\textquotedblleft{}Guilty; owing\textquotedblright{}.) Say it once more.
\end{tcolorbox}
\section[unschuldig]{unschuldig — innocent}

\label{lesson:GE-C286-unschuldig}



\begin{quote}
Before the new one: say the German for to sentence, then the German for guilty.
\end{quote}

\subsection*{You'll want to know: unschuldig}
\textbf{unschuldig} — \textquotedblleft{}innocent\textquotedblright{}.

\begin{quote}
Am Ende war die Frau unschuldig.
\end{quote}

\begin{quote}
\emph{In the end the woman was innocent.}
\end{quote}

The opposite of \emph{schuldig}.

\begin{cousinweb}
\emph{unschuldig} joins \emph{un-} \textquotedblleft{}not\textquotedblright{} and \emph{schuldig} \textquotedblleft{}guilty\textquotedblright{}, from Old High German \emph{sculdig} (EtymDB).
\end{cousinweb}

\begin{grammarlens}[title={What you've built}]
One more word for telling, explaining and describing at B1.
\end{grammarlens}

\subsection*{Your turn}
\begin{itemize}
\raggedright
  \item \emph{Say it:} \emph{unschuldig}
  \item \emph{Say it:} \emph{unschuldig}, once more
  \item \emph{Say it:} a sentence of your own with \emph{unschuldig}
  \item \emph{Recall:} say the German for to sentence, then the German for guilty, then say \emph{unschuldig} again
\end{itemize}

\begin{tcolorbox}[breakable,colback=teal!4,colframe=teal!35!black,title={Before you move on}]
What does unschuldig mean? (\textquotedblleft{}Innocent\textquotedblright{}.) Say it once more.
\end{tcolorbox}
